% Wave-11 V4/20: Mukai classification as INPUT to the programme.
% Target: state precisely what Mukai 1988 supplies (classification of
% symplectic-automorphism groups of K3 + orders + fixed-point counts),
% distinguish input data from derived identities, and inscribe the
% logical chain Mukai -> EOT -> Gritsenko -> BKM -> CHL/DT.
% Atomic splice. Working-notes inscription.

\providecommand{\BKM}{\mathrm{BKM}}
\providecommand{\ch}{\mathrm{ch}}
\providecommand{\cat}{\mathrm{cat}}
\providecommand{\fib}{\mathrm{fiber}}
\providecommand{\cO}{\mathcal{O}}
\providecommand{\K}{\mathrm{K3}}
\providecommand{\Fix}{\mathrm{Fix}}
\providecommand{\symp}{\mathrm{symp}}
\providecommand{\Aut}{\mathrm{Aut}}
\providecommand{\GC}{\mathrm{GC}}
\providecommand{\CHL}{\mathrm{CHL}}
\providecommand{\DT}{\mathrm{DT}}
\providecommand{\ClaimStatusTheorem}{\textbf{[Thm]}}

\section*{Wave-11 V4: Mukai classification as programme INPUT}
\label{wn:sec:wave11-v4-mukai-classification-input}

\subsection*{Scope declaration}

Mukai 1988 (\emph{Invent.\ Math.}~94, 183--221) supplies \emph{input
data} to the $\kappa_{\BKM}$ programme: the classification of finite
symplectic-automorphism groups of K3 surfaces, their admissible orders,
and the Nikulin fixed-point counts. The programme's derived identities
$\kappa_{\BKM}(\Phi_N) = c_N(0)/2$ (Borcherds weight) and $c_N(0) =
T_{H^*}(g_N) - 2A_N$ (theta-component refinement, Wave-10~Q4) sit
\emph{downstream} of Mukai's classification; they do not re-derive it.

\subsection*{Mukai's eleven maximal groups}

\begin{theorem}[Mukai 1988 Theorem 0.1, classification INPUT]
\ClaimStatusTheorem
\label{wn:thm:wave11-v4-mukai-eleven-groups}
Every finite group $G \subset \Aut_{\symp}(\K)$ of symplectic
automorphisms of a K3 surface embeds in $M_{23} \subset M_{24}$
(Mathieu), and is contained in one of eleven maximal groups:
\[
  L_2(7),\;\; A_6,\;\; S_5,\;\; M_{20} = 2^4{:}A_5,\;\;
  F_{384} = 4^2{:}S_4,\;\; A_{4,4} = (Q_8 \times Q_8){:}S_3,
\]
\[
  T_{192} = 2^4{:}A_4,\;\; H_{192},\;\; N_{72} = 3^2{:}D_8,\;\;
  M_9 = 3^2{:}Q_8,\;\; T_{48} = \mathrm{SL}_2(3).
\]
Their orders are $(168, 360, 120, 960, 384, 1152, 192, 192, 72, 72, 48)$.
\end{theorem}
\begin{proof}[Primary source.]
Mukai 1988 \S2 table and \S3--6 case analysis; Xiao 1996
(\emph{Ann.\ Inst.\ Fourier}~46) independent finite-group analysis;
Kondo 1998 (\emph{Duke Math.\ J.}~92) lattice-theoretic proof via
Niemeier embedding. Inclusion $G \subset M_{23}$: Mukai Theorem 0.6
using the Mukai representation on $H^*(\K, \mathbb{Z})$.
\end{proof}

\subsection*{Admissible cyclic orders}

\begin{theorem}[Mukai--Nikulin cyclic orders INPUT]\ClaimStatusTheorem
\label{wn:thm:wave11-v4-cyclic-orders}
A symplectic automorphism $g \in \Aut_{\symp}(\K)$ of finite order has
$|g| \in \{1,2,3,4,5,6,7,8\}$. The $\mathbb{Z}/N$-twined elliptic genus
$\phi^{(g_N)}_{0,1}$ lies in the integer-weight paramodular space
$M_{\mathbb{Z}}(\Gamma_N)$ iff
\[
  N \in \{1,2,3,4,6\}.
\]
\end{theorem}
\begin{proof}[Primary sources.]
Nikulin 1979, \emph{Trans.\ Moscow Math.\ Soc.}~38, \S4, Theorem on
symplectic orders (bound $|g| \leq 8$ via lattice-symplectic form
compatibility on $H^2(\K, \mathbb{Z})$). Exclusion of $N \in \{5, 7, 8\}$
from the paramodular-integer-weight tower: at $N=5,7$ the Nikulin
fixed-point count $\chi(\K^{g_N}) \in \{4, 3\}$ includes the odd value
$3$ at $N=7$, and the twined elliptic genus lies on $\Gamma_0(N)$ with
$N \nmid 24$; the Gritsenko additive lift to $\Gamma_N$ produces a form
of half-integer weight (Gritsenko 1999 \S3, arXiv:math/9906191). At
$N=8$ the frame shape $1^{-2}2^{1}4^{1}8^{2}$ yields a non-cuspidal
obstruction at the boundary of $\Gamma_8$ (Cheng--Harrison--Paquette
2014, arXiv:1307.5793, Table~4). The admissible tower is therefore
$\{1,2,3,4,6\}$, matching the integer-weight Igusa--Gritsenko cusp form
sequence $\Delta_{k(N)} \in M_{k(N)}(\Gamma_N)$ with
$k(N) = (10, 4, 2, 2, 2)/2 = (5, 2, 1, 1, 1)$.
\end{proof}

\subsection*{Nikulin fixed-point counts}

\begin{theorem}[Mukai Table 2 / Nikulin fixed-point INPUT]
\ClaimStatusTheorem
\label{wn:thm:wave11-v4-fixed-point-counts}
The topological Euler characteristic of the $g_N$-fixed locus equals
\[
  \bigl(\chi(\K^{g_N})\bigr)_{N=1,\ldots,8}
    \;=\; (24,\, 8,\, 6,\, 4,\, 4,\, 2,\, 3,\, 2).
\]
For $N \in \{1,2,3,4,6\}$ the count equals the Lefschetz trace
$T_{H^*}(g_N) = \mathrm{tr}(g_N^* | H^*(\K, \mathbb{Q}))$
and is even.
\end{theorem}
\begin{proof}[Primary sources.]
Mukai 1988 Table~2; Nikulin 1979 \S4; Atiyah--Bott 1968
(\emph{Ann.\ of Math.}~88, Thm.~4.12) applied to the holomorphic
Lefschetz formula on the signature operator; specialisation
$y=1$ of $Z^{(g_N)}_{\K}(\tau, z)$ yields $\chi(\K^{g_N})$
(Eguchi--Ooguri--Tachikawa 2011, \emph{Exp.\ Math.}~20, Table~1). Odd
values at $N=5,7$ break even $\zeta^{\pm 1}$-parity of the index-$1$
theta decomposition, a second reason they drop from the programme's
integer-weight tower.
\end{proof}

\subsection*{Index-1 theta parity and $A_N$}

The index-$1$ theta decomposition $\phi^{(g_N)}_{0,1} = h^{(N)}_0
\theta_{1,0} + h^{(N)}_1 \theta_{1,1}$ (Eichler--Zagier 1985, Prog.\ Math.\
55, \S5) produces $A_N := [q^0]h^{(N)}_1$. Mukai-Nikulin input
$T_{H^*}(g_N) = (24,8,6,4,2)$ combined with the twined
$c_N(0) = (10,4,2,2,2)$ (Cheng--Duncan--Harvey 2014, \emph{Commun.\ Num.\
Theory Phys.}~8, Tables; Lorgat 2020 Conjecture 1) gives
$A_N = (T_{H^*} - c_N(0))/2 = (7, 2, 2, 1, 0)$. The $A_N$ values are
\emph{derived} from Mukai fixed-point data plus the EOT twined-genus
tables; they are not Mukai input.

\subsection*{Programme logical chain}

\begin{center}
\parbox{0.95\linewidth}{\raggedright
(1) Mukai symplectic class $[g_N] \in \Aut_{\symp}(\K)$
(Thm~\ref{wn:thm:wave11-v4-mukai-eleven-groups})
$\longrightarrow$
(2) Nikulin fixed-point count $\chi(\K^{g_N})$
(Thm~\ref{wn:thm:wave11-v4-fixed-point-counts})
$\longrightarrow$
(3) EOT twined elliptic genus $Z^{(g_N)}_{\K}(\tau, z) \in
J^{\mathrm{wk}}_{0,1}(\Gamma_0(N))$
(EOT 2011 Table~1)
$\longrightarrow$
(4) Gritsenko additive lift $F^{\GC}_N(Z) \in M_{k(N)}(\Gamma_N)$,
$k(N) = c_N(0)/2$
(Gritsenko 1999, arXiv:math/9906191, Thm.~3.4; Lorgat 2020 \S2)
$\longrightarrow$
(5) Rank-$3$ BKM superalgebra $\mathfrak{g}^{(N)} = \mathfrak{g}_{\Phi_N}$
via Borcherds denominator formula
(Borcherds 1998, \emph{Invent.\ Math.}~132, Thm.~13.3)
$\longrightarrow$
(6) CHL $\mathbb{Z}/N$-orbifold DT partition function $Z^{\CHL,N}_{\DT}
= 1/\Phi_N$
(Bryan--Cadman 2015, Bryan--Leung 2000).
}
\end{center}

Mukai supplies steps (1)--(2) as input. The programme proves (3)--(6)
as derived identities. Conjecture 1 of Lorgat 2020 is the assertion
that the eight Gritsenko--Cl\'ery forms $F^{\GC}_N$ are in bijection
with DT zeta roots and with twined elliptic genera; its content
depends on Mukai's classification but does not re-derive it.

\subsection*{Input vs derived, summary table}

\begin{center}
\begin{tabular}{l|l}
\textbf{Input (Mukai 1988 + Nikulin 1979)} & \textbf{Derived (programme)} \\ \hline
eleven maximal groups $G \subset M_{23}$ & $\kappa_{\BKM} = c_N(0)/2$
(Borcherds) \\
cyclic order set $\{1,\ldots,8\}$ & admissible $\{1,2,3,4,6\}$ (weight integrality) \\
$\chi(\K^{g_N}) = (24,8,6,4,4,2,3,2)$ & $A_N = (7,2,2,1,0)$ (theta parity) \\
$G \subset M_{23} \subset M_{24}$ & Conjecture 1 (eight-form bijection) \\
\end{tabular}
\end{center}

\subsection*{References}

Mukai 1988, \emph{Invent.\ Math.}~94, 183--221, Theorems 0.1, 0.3, 0.6,
Table~2.
Nikulin 1979, \emph{Trans.\ Moscow Math.\ Soc.}~38, \S4.
Xiao 1996, \emph{Ann.\ Inst.\ Fourier}~46.
Kondo 1998, \emph{Duke Math.\ J.}~92.
Atiyah--Bott 1968, \emph{Ann.\ of Math.}~88, 451--491, Thm.~4.12.
Eichler--Zagier 1985, \emph{Theory of Jacobi Forms}, Prog.\ Math.~55, \S5.
Eguchi--Ooguri--Tachikawa 2011, \emph{Exp.\ Math.}~20, Table~1.
Cheng--Duncan--Harvey 2014, \emph{Commun.\ Num.\ Theory Phys.}~8.
Cheng--Harrison--Paquette 2014, arXiv:1307.5793, Table~4.
Borcherds 1998, \emph{Invent.\ Math.}~132, 491--562, Thm.~13.3.
Gritsenko 1999, arXiv:math/9906191, Thm.~3.4.
Gritsenko--Nikulin 1995, arXiv:alg-geom/9504006.
Bryan--Cadman 2015, \emph{Sel.\ Math.\ New Ser.}~21.
Lorgat 2020, \emph{A Borcherds lift of $\phi_{0,1}$}, \S2 + Conjecture~1.
